\section{Introduction} \label{sec:introduction}
In the case of Bitcoin, hash-based signature schemes are the most compatible among all quantum-resistant signature algorithms. Still, the smallest stateless hash-based signature proposed by NIST (SLH-DSA-128s) is \emph{7,856} bytes \cite{nist}, which is more than $100\times$ larger than the current Schnorr signature. The question is whether we can have something smaller and more practical without significant deviation from the standard.

SHRINCS \cite{delving-post} proposes a hybrid signature consisting of a stateful and a stateless part, with signature size ranging from \emph{324} bytes (stateful) to \emph{X} bytes (stateless). This document attempts to answer the question of what the value of \emph{X} might be while keeping other signature metrics (KeyGen, SigGen, SigVer, etc.) practical.

In this document, we present the methodology for searching parameters of SLH-DSA-compatible signatures: how changes in the hyper-tree structure and other parameters affect the signature metrics. As mentioned above, we deliberately stay within very conservative bounds and avoid any changes to the signature sub-algorithms in order to keep the construction as compatible as possible with the SLH-DSA standard.

\textit{Along the way, we will occasionally reference promising proposals such as hyper-tree pruning \cite{conduition-blog} and the +C compression technique \cite{mike-paper}, which allow for even better shrinking results. Although some of them can be considered partially compatible, we do not rely on them in the evaluation framework. We still think it is useful to outline the potential of such techniques and to let the reader tweak the parameters accordingly.}

The main contributions of this report are (1) a sweep script, (2) the filtering pipeline, and (3) a recommended candidate tuple.